
\documentclass[11pt,a4paper]{article}
\usepackage[utf8]{inputenc}
\usepackage[T1]{fontenc}
\usepackage[french]{babel}
\usepackage{lmodern}
\usepackage{microtype}
\usepackage[a4paper,margin=1.65cm,headheight=15pt]{geometry}
\usepackage{amsmath,amssymb,mathtools}
\usepackage{array,tabularx,multicol,booktabs}
\usepackage{enumitem}
\usepackage[most]{tcolorbox}
\usepackage{xcolor}
\usepackage{tikz}
\usetikzlibrary{arrows.meta,calc,positioning,angles,quotes}
\usepackage{pdflscape}
\usepackage{fancyhdr}
\usepackage{hyperref}
\usepackage{needspace}
\usepackage{pagecolor}

\definecolor{fondpage}{RGB}{247,248,250}
\definecolor{encre}{RGB}{31,35,40}
\definecolor{bleu}{RGB}{40,91,135}
\definecolor{vert}{RGB}{55,125,92}
\definecolor{orange}{RGB}{178,105,42}
\definecolor{violet}{RGB}{101,77,145}
\definecolor{rouge}{RGB}{170,72,72}
\definecolor{gris}{RGB}{86,91,99}
\definecolor{bleufonce}{RGB}{28,55,120}
\definecolor{bleuclair}{RGB}{238,243,255}
\definecolor{grisclair}{RGB}{245,245,245}
\definecolor{tableline}{RGB}{45,45,45}
\definecolor{softgray}{RGB}{246,247,249}
\definecolor{deepblue}{RGB}{25,62,115}
\definecolor{accent}{RGB}{20,82,130}

\hypersetup{colorlinks=true,linkcolor=bleufonce,urlcolor=bleufonce,
pdftitle={Parcours de revision -- Produit scalaire -- Premiere specialite},pdfauthor={}}

\setlength{\parindent}{0pt}
\setlength{\parskip}{0.25em}
\renewcommand{\arraystretch}{1.13}
\setlength{\tabcolsep}{4pt}
\setlist[itemize]{leftmargin=1.25em,itemsep=0.15em,topsep=0.2em}
\setlist[enumerate,1]{label=\textbf{\arabic*.}, leftmargin=1.05cm, itemsep=0.35em, topsep=0.2em}
\setlist[enumerate,2]{label=\textbf{\alph*.}, leftmargin=0.85cm, itemsep=0.25em, topsep=0.15em}

\pagestyle{fancy}
\fancyhf{}
\cfoot{\thepage}

\newcommand{\R}{\mathbb{R}}
\newcommand{\N}{\mathbb{N}}
\newcommand{\vect}[1]{\overrightarrow{#1}}
\newcommand{\norme}[1]{\left\lVert #1\right\rVert}
\newcommand{\sourceq}[1]{ {\scriptsize\textcolor{bleufonce}{(site Premiere q#1)}}}

\newcounter{question}
\newcounter{reponse}

\newenvironment{blocquestions}[1]{%
  \begin{tcolorbox}[enhanced,breakable,colback=bleuclair,colframe=bleufonce,boxrule=0.6pt,arc=2mm,left=2mm,right=2mm,top=2mm,bottom=1.5mm,title={\bfseries #1},coltitle=white,colbacktitle=bleufonce,before skip=3mm,after skip=3mm, fontupper=\large]
  \large
  \begin{enumerate}[leftmargin=*,label=\textbf{\arabic*.},itemsep=0.82em,topsep=0.5em]
  \setcounter{enumi}{\value{question}}
}{%
  \setcounter{question}{\value{enumi}}
  \end{enumerate}
  \end{tcolorbox}
}

\newenvironment{blocreponses}[1]{%
  \begin{tcolorbox}[enhanced,breakable,colback=bleuclair,colframe=bleufonce,boxrule=0.6pt,arc=2mm,left=2mm,right=2mm,top=2mm,bottom=1.5mm,title={\bfseries #1},coltitle=white,colbacktitle=bleufonce,before skip=3mm,after skip=3mm, fontupper=\large]
  \large
  \begin{enumerate}[leftmargin=*,label=\textbf{\arabic*.},itemsep=0.42em,topsep=0.5em]
  \setcounter{enumi}{\value{reponse}}
}{%
  \setcounter{reponse}{\value{enumi}}
  \end{enumerate}
  \end{tcolorbox}
}

\newtcolorbox{correctionbox}[1]{enhanced,breakable,colback=grisclair,colframe=black!55,boxrule=0.55pt,arc=2mm,left=2mm,right=2mm,top=2mm,bottom=1.5mm,title=\textbf{#1},coltitle=black,colbacktitle=black!12,before skip=2mm,after skip=2mm}
\newtcolorbox{methodebox}[1]{enhanced,breakable,colback=white,colframe=bleu!70!black,boxrule=0.55pt,arc=2mm,left=2mm,right=2mm,top=2mm,bottom=1.5mm,title=\textbf{#1},coltitle=white,colbacktitle=bleu!70!black,before skip=2mm,after skip=2mm}

\newcommand{\titreSujet}[2]{%
  \subsection*{#1}
  \addcontentsline{toc}{subsection}{#1}
  \textit{Référence / inspiration : #2}\par\medskip\hrule\medskip
}
\newcommand{\qcm}[4]{%
\begin{center}\begin{tabularx}{0.96\linewidth}{@{}XX@{}}
\textbf{a.} #1 & \textbf{b.} #2\\[1mm]
\textbf{c.} #3 & \textbf{d.} #4
\end{tabularx}\end{center}}

\tcbset{enhanced,arc=2mm,outer arc=2mm,boxrule=0.42pt,left=1.15mm,right=1.15mm,top=0.75mm,bottom=0.75mm,boxsep=0.35mm,before skip=0.8mm,after skip=0.8mm,fonttitle=\bfseries\footnotesize,coltitle=encre,before upper=\raggedright\fontsize{7.05}{7.55}\selectfont}
\newtcolorbox{fichebox}[3][]{colback=white,colframe=#2!72!black,colbacktitle=#2!18,coltitle=#2!50!black,borderline west={0.9mm}{0pt}{#2!78!black},title={#3},drop fuzzy shadow=black!5,#1}
\newcommand{\tagcol}[2]{\begin{tcolorbox}[enhanced,colback=#1!11,colframe=#1!45!black,boxrule=0.42pt,arc=2mm,left=1mm,right=1mm,top=0.45mm,bottom=0.45mm,before skip=0mm,after skip=0.85mm,fontupper=\bfseries\scriptsize,colupper=#1!55!black]\centering #2\end{tcolorbox}}
\newtcolorbox{panel}[1]{colback=fondpage,colframe=fondpage,boxrule=0pt,arc=0mm,left=0mm,right=0mm,top=0mm,bottom=0mm,height=168mm,valign=top,before skip=0mm,after skip=0mm}

% Mini-schémas TikZ inspirés des fiches visuelles de Première
\newcommand{\schemaAngleProduit}{%
\begin{center}
\begin{tikzpicture}[scale=0.48,>=Latex,line cap=round]
  \coordinate (A) at (0,0);
  \coordinate (B) at (2.45,0);
  \coordinate (C) at (1.55,1.35);
  \draw[->,very thick,bleu!85!black] (A)--(B) node[midway,below=2pt] {\scriptsize $\vect{AB}$};
  \draw[->,very thick,orange!88!black] (A)--(C) node[midway,above left=1pt] {\scriptsize $\vect{AC}$};
  \pic[draw=gris,thick,angle radius=0.46cm,"{\scriptsize $\theta$}",angle eccentricity=1.45] {angle=B--A--C};
  \fill (A) circle (1.2pt) node[below left=1pt] {\scriptsize $A$};
  \fill (B) circle (1.2pt) node[below right=1pt] {\scriptsize $B$};
  \fill (C) circle (1.2pt) node[above=1pt] {\scriptsize $C$};
\end{tikzpicture}
\end{center}}
\newcommand{\schemaProjectionProduit}{%
\begin{center}
\begin{tikzpicture}[scale=0.50,>=Latex,line cap=round]
  \coordinate (A) at (0,0);
  \coordinate (B) at (3.65,0);
  \coordinate (H) at (2.35,0);
  \coordinate (C) at (2.35,1.55);

  % base et vecteurs
  \draw[->,very thick,bleu!85!black] (A)--(B) node[midway,below=6pt] {\scriptsize $\vect{AB}$};
  \draw[->,very thick,orange!88!black] (A)--(C) node[midway,above left=1pt] {\scriptsize $\vect{AC}$};

  % projection
  \draw[dashed,gris,thick] (C)--(H);
  \draw[gris,thick] ($(H)+(-0.18,0)$)--++(0,0.18)--++(0.18,0);

  % points
  \fill (A) circle (1.2pt) node[below left=1pt] {\scriptsize $A$};
  \fill (B) circle (1.2pt) node[below right=1pt] {\scriptsize $B$};
  \fill (C) circle (1.2pt) node[above=1pt] {\scriptsize $C$};
  \fill (H) circle (1.2pt) node[below=2pt] {\scriptsize $H$};
\end{tikzpicture}
\end{center}}
\newcommand{\schemaNormalDroite}{%
\begin{center}
\begin{tikzpicture}[scale=0.42,>=Latex,line cap=round]
  % Droite : segment simple, sans flèche.
  \draw[thick,bleu!80!black,line cap=butt] (-1.25,-.48)--(2.25,.86);
  \node[bleu!80!black,right=2pt] at (2.25,.86) {\scriptsize $d$};
  \coordinate (P) at (.5,.25);
  \draw[->,thick,rouge!85!black] (P)--++(-.55,1.45) node[above] {\scriptsize $\vec n$};
\end{tikzpicture}
\end{center}}
\newcommand{\schemaCercleDiametre}{%
\begin{center}
\begin{tikzpicture}[scale=0.48,line cap=round]
  \coordinate (A) at (-1.60,0);
  \coordinate (B) at (1.60,0);
  \coordinate (O) at (0,0);
  \coordinate (M) at (0.42,1.544);

  % Cercle de diamètre [AB]
  \draw[thick,bleu!80!black] (O) circle (1.60);

  % Diamètre et triangle inscrit
  \draw[thick] (A)--(B);
  \draw[thick] (A)--(M)--(B);

  % Angle droit en M
  \pic[draw=rouge!70!black,thick,angle radius=.28cm] {right angle=A--M--B};

  % Points et labels
  \fill (A) circle (1.15pt) node[below left=2pt] {\scriptsize $A$};
  \fill (B) circle (1.15pt) node[below right=2pt] {\scriptsize $B$};
  \fill (M) circle (1.15pt) node[above=2pt] {\scriptsize $M$};
\end{tikzpicture}
\end{center}}
\newcommand{\schemaTangenteCercle}{%
\begin{center}
\begin{tikzpicture}[scale=0.42,>=Latex,line cap=round]
  \coordinate (I) at (0,0);\coordinate (T) at (1.45,.55);
  \draw[thick,bleu!80!black] (I) circle (1.55);
  \draw[thick,rouge!85!black] ($(T)+(-.55,1.45)$)--($(T)+(.55,-1.45)$) node[below] {\scriptsize tangente};
  \draw[->,thick,orange!90!black] (I)--(T) node[midway,below] {\scriptsize rayon};
  \fill (I) circle (1.2pt) node[left] {\scriptsize $I$};\fill (T) circle (1.2pt) node[right] {\scriptsize $T$};
  \draw[gris] (T) ++(-.07,.17)--++(.17,.07)--++(.07,-.17);
\end{tikzpicture}
\end{center}}

\newcommand{\titrepage}{\begin{tcolorbox}[colback=white,colframe=bleu!70!black,boxrule=0.65pt,arc=3mm,left=2.5mm,right=2.5mm,top=1.15mm,bottom=1.15mm,before skip=0mm,after skip=1.4mm,drop fuzzy shadow=black!10]
{\Large\bfseries Parcours de révision -- Produit scalaire}\hfill {\normalsize Première spécialité -- sans calculatrice}
\end{tcolorbox}}

\newcommand{\ficheRecap}{%
\begin{landscape}
\pagecolor{fondpage}
\thispagestyle{empty}
\titrepage
\begin{tabularx}{\linewidth}{@{}X@{\hspace{1.6mm}}X@{\hspace{1.6mm}}X@{\hspace{1.6mm}}X@{}}

\begin{panel}{}
\tagcol{gris}{Rappels de Seconde}
\begin{fichebox}{gris}{Coordonnées de $\vect{AB}$}
Si $A(x_A;y_A)$ et $B(x_B;y_B)$, alors
\[
\vect{AB}=\begin{pmatrix}x_B-x_A\\ y_B-y_A\end{pmatrix}.
\]
\end{fichebox}
\begin{fichebox}{gris}{Distance}
Dans un repère orthonormé :
\[
AB=\sqrt{(x_B-x_A)^2+(y_B-y_A)^2}.
\]
Donc $AB^2=(x_B-x_A)^2+(y_B-y_A)^2$.
\end{fichebox}
\begin{fichebox}{gris}{Milieu d'un segment}
Le milieu $I$ de $[AB]$ a pour coordonnées :
\[
I\left(\frac{x_A+x_B}{2};\frac{y_A+y_B}{2}\right).
\]
\end{fichebox}
\begin{fichebox}{gris}{À utiliser avant le produit scalaire}
Pour calculer $\vect{AB}\cdot\vect{AC}$, on commence souvent par calculer les coordonnées de $\vect{AB}$ et de $\vect{AC}$.
\end{fichebox}
\end{panel}
&
\begin{panel}{}
\tagcol{bleu}{Définitions et calculs}
\begin{fichebox}{bleu}{Produit scalaire en coordonnées}
Dans un repère orthonormé, si
\[
\vec u=\begin{pmatrix}x\\y\end{pmatrix},\qquad
\vec v=\begin{pmatrix}x'\\y'\end{pmatrix},
\]
alors
\[
\vec u\cdot\vec v=xx'+yy'.
\]
\end{fichebox}
\begin{fichebox}{bleu}{Avec les longueurs}
\schemaAngleProduit
\[
\vect{AB}\cdot\vect{AC}=AB\times AC\times\cos(\widehat{BAC}).
\]
En particulier : \(\vec u\cdot\vec u=\norme{\vec u}^{2}\).
\end{fichebox}
\begin{fichebox}{bleu}{Propriétés de calcul}
Symétrie : \(\vec u\cdot\vec v=\vec v\cdot\vec u\).

Bilinéarité :
\[
(\vec u+\vec v)\cdot\vec w=\vec u\cdot\vec w+\vec v\cdot\vec w,
\]
\[
(k\vec u)\cdot\vec v=k(\vec u\cdot\vec v).
\]
Identité utile :
\[
\norme{\vec u+\vec v}^{2}=\norme{\vec u}^{2}+2\vec u\cdot\vec v+\norme{\vec v}^{2}.
\]
\end{fichebox}
\begin{fichebox}{bleu}{Al-Kashi}
Dans un triangle \(ABC\) :
\[
BC^2=AB^2+AC^2-2\,AB\times AC\cos(\widehat{BAC}).
\]
\end{fichebox}
\end{panel}
&
\begin{panel}{}
\tagcol{vert}{Orthogonalité et angle}
\begin{fichebox}{vert}{Orthogonalité}
Deux vecteurs non nuls \(\vec u\) et \(\vec v\) sont orthogonaux si et seulement si :
\[
\vec u\cdot\vec v=0.
\]
Donc deux droites sont perpendiculaires si leurs vecteurs directeurs ont un produit scalaire nul.
\end{fichebox}
\begin{fichebox}{vert}{Calculer un angle}
Si \(\vec u\) et \(\vec v\) sont non nuls :
\[
\cos\theta=\frac{\vec u\cdot\vec v}{\norme{\vec u}\,\norme{\vec v}}.
\]
Le signe du produit scalaire renseigne sur l’angle : aigu si positif, droit si nul, obtus si négatif.
\end{fichebox}
\begin{fichebox}{vert}{Projeté orthogonal}
\schemaProjectionProduit
Si \(H\) est le projeté orthogonal de \(C\) sur \((AB)\), alors :
\[
\vect{AB}\cdot\vect{AC}=\vect{AB}\cdot\vect{AH}.
\]
\end{fichebox}
\end{panel}
&
\begin{panel}{}
\tagcol{orange}{Applications repérées}
\begin{fichebox}{orange}{Vecteur normal}
\schemaNormalDroite
Pour une droite \(ax+by+c=0\), un vecteur normal est
\[
\vec n=\begin{pmatrix}a\\b\end{pmatrix}.
\]
Un vecteur directeur possible est \(\vec d=\begin{pmatrix}-b\\a\end{pmatrix}\).
\end{fichebox}
\begin{fichebox}{orange}{Équation de cercle}
Le cercle de centre \(I(x_I;y_I)\) et de rayon \(r\) a pour équation :
\[
(x-x_I)^2+(y-y_I)^2=r^2.
\]
\end{fichebox}
\begin{fichebox}{orange}{Cercle de diamètre}
\schemaCercleDiametre
Un point \(M\) appartient au cercle de diamètre \([AB]\) si et seulement si :
\[
\vect{MA}\cdot\vect{MB}=0.
\]
\end{fichebox}
\end{panel}
\end{tabularx}
\clearpage
\nopagecolor
\end{landscape}}

\begin{document}
\begin{center}
{\LARGE\bfseries Corrigé -- Parcours de révision -- Produit scalaire}\\[1mm]
{\large Première générale -- spécialité mathématiques}\\[1mm]
{\normalsize Sans calculatrice}
\end{center}
\tableofcontents\newpage
\phantomsection\addcontentsline{toc}{section}{Fiche récapitulative}
\ficheRecap
\section{Correction des questions flash}
\setcounter{reponse}{0}
\begin{blocreponses}{Questions flash 1 à 12}
\item $2\times 3+(-1)\times4=2$.
\item $1\times4+2\times(-2)=0$. Oui, ils sont orthogonaux.
\item $\overrightarrow{AB}=(5-1;-1-3)=(4;-4)$.
\item $\overrightarrow{AB}=(3;1)$, $\overrightarrow{AC}=(2;-6)$ donc $3\times2+1\times(-6)=0$.
\item Ils sont orthogonaux.
\item $\norme{\vec u}=\sqrt{6^2+8^2}=10$.
\item On note $\theta$ l'angle formé par $\vec u$ et $\vec v$. Alors $\vec u\cdot\vec v=\norme{\vec u}\,\norme{\vec v}\cos\theta$. Donc $6=2\times3\times\cos\theta$, d'où $\cos\theta=1$.
\item C’est un angle droit.
\item Un vecteur normal est $\vec n(3;-2)$.
\item Un vecteur directeur possible est $\vec d(2;3)$.
\item $2\times1+2-4=0$, donc oui.
\item $(x-2)^2+(y+1)^2=9$.
\end{blocreponses}
\begin{blocreponses}{Questions flash 13 à 24}
\item Oui : $(5-2)^2+(-1+1)^2=9$.
\item Centre $I(-1;4)$ et rayon $5$.
\item $\overrightarrow{MA}\cdot\overrightarrow{MB}=0$.
\item Le vecteur $\overrightarrow{IT}$ est normal à la tangente.
\item $\vec u\cdot\vec u=25$.
\item $2\times 3+m\times6=0$, donc $m=-1$.
\item $1\times8+4\times(-2)=0$, donc oui.
\item $AB=\sqrt{4^2+3^2}=5$.
\item $\overrightarrow{AB}=(3;0)$, $\overrightarrow{AC}=(0;4)$, donc le produit vaut $0$.
\item L’angle entre les vecteurs est aigu.
\item L’angle entre les vecteurs est obtus.
\item Un vecteur normal est $\vec n(1;4)$.
\end{blocreponses}
\begin{blocreponses}{Questions flash 25 à 36}
\item $2(x-1)+5(y+1)=0$, soit $2x+5y+3=0$.
\item Un vecteur directeur possible est $(3;1)$.
\item $\overrightarrow{AB}=(4;2)$, $\overrightarrow{AC}=(1;-2)$, produit $=4-4=0$, angle droit en $A$.
\item Centre $(2;0)$, rayon $2$ : $(x-2)^2+y^2=4$.
\item Oui, car $r^2=10$ donc $r=\sqrt{10}$.
\item L’angle vaut $0^\circ$ : vecteurs de même sens.
\item L’angle vaut $180^\circ$ : vecteurs de sens opposés.
\item $\overrightarrow{BA}=(1;0)$, $\overrightarrow{BC}=(0;2)$ donc produit $0$.
\item $\overrightarrow{MA}\cdot\overrightarrow{MB}=0$.
\item On montre que le produit scalaire des deux vecteurs portés par les côtés de l’angle vaut $0$.
\item C’est un vecteur normal.
\item $4\times5\times \frac12=10$.
\end{blocreponses}

\section{Correction du QCM}
\begin{enumerate}
\item Réponse c : $2\times(-1)+3\times4=10$.
\item Réponse b : $5\times1+1\times(-5)=0$.
\item Réponse b : un vecteur normal est $(-2;3)$.
\item Réponse b : $(x-1)^2+(y-2)^2=16$.
\item Réponse a : critère du cercle de diamètre.
\item Réponse c : le produit scalaire négatif indique un angle obtus.
\item Réponse b : la tangente est perpendiculaire au rayon au point de tangence.
\item Réponse a : $AB=\sqrt{3^2+4^2}=5$.
\end{enumerate}

\section{Correction des sujets type bac / E3C}

\subsection*{Sujet 1 -- Triangle, angle et orthogonalité}
\begin{correctionbox}{1. Coordonnées de vecteurs}
\[
\overrightarrow{AB}=(5-1;4-2)=(4;2),\qquad
\overrightarrow{AC}=(3-1;-2-2)=(2;-4).
\]
\end{correctionbox}
\begin{correctionbox}{2. Produit scalaire}
\[
\overrightarrow{AB}\cdot\overrightarrow{AC}=4\times2+2\times(-4)=8-8=0.
\]
\end{correctionbox}
\begin{correctionbox}{3. Orthogonalité}
Le produit scalaire est nul, donc les vecteurs $\overrightarrow{AB}$ et $\overrightarrow{AC}$ sont orthogonaux. Les droites $(AB)$ et $(AC)$ sont donc perpendiculaires.
\end{correctionbox}
\begin{correctionbox}{4. Longueurs}
\[
AB=\sqrt{4^2+2^2}=\sqrt{20}=2\sqrt5,
\qquad
AC=\sqrt{2^2+(-4)^2}=\sqrt{20}=2\sqrt5.
\]
\end{correctionbox}
\begin{correctionbox}{5. Cosinus de l’angle}
\[
\cos(\widehat{BAC})=\frac{\overrightarrow{AB}\cdot\overrightarrow{AC}}{AB\times AC}=\frac{0}{(2\sqrt5)(2\sqrt5)}=0.
\]
\end{correctionbox}
\begin{correctionbox}{6. Nature du triangle}
Comme $(AB)\perp(AC)$, le triangle $ABC$ est rectangle en $A$.
\end{correctionbox}
\begin{correctionbox}{7. Produit scalaire en $B$}
\[
\overrightarrow{BA}=(-4;-2),\qquad \overrightarrow{BC}=(-2;-6).
\]
Donc :
\[
\overrightarrow{BA}\cdot\overrightarrow{BC}=(-4)(-2)+(-2)(-6)=8+12=20.
\]
Ce produit scalaire est positif : l’angle $\widehat{ABC}$ est aigu.
\end{correctionbox}

\subsection*{Sujet 2 -- Droite, vecteur normal et projeté orthogonal}
\begin{correctionbox}{1. Vecteur normal}
Pour une droite $ax+by+c=0$, un vecteur normal est $(a;b)$. Ici :
\[
\vec n=(2;-1).
\]
\end{correctionbox}
\begin{correctionbox}{2. Droite perpendiculaire}
La droite $\Delta$ est perpendiculaire à $d$. Comme $\vec n$ est normal à $d$, il est directeur d’une droite perpendiculaire à $d$.
\end{correctionbox}
\begin{correctionbox}{3. Équation paramétrique de $\Delta$}
La droite $\Delta$ passe par $A(4;1)$ et a pour vecteur directeur $(2;-1)$, donc :
\[
\left\{\begin{array}{rcl}
x&=&4+2t\\
y&=&1-t
\end{array}\right.
\qquad t\in\mathbb R.
\]
\end{correctionbox}
\begin{correctionbox}{4. Coordonnées du point $H$}
On remplace dans l’équation de $d$ :
\[
2(4+2t)-(1-t)-3=0.
\]
Donc :
\[
8+4t-1+t-3=0,
\qquad 4+5t=0,
\qquad t=-\frac45.
\]
Ainsi :
\[
H\left(4+2\times\left(-\frac45\right);1-\left(-\frac45\right)\right)
=\left(\frac{12}{5};\frac95\right).
\]
\end{correctionbox}
\begin{correctionbox}{5. Projeté orthogonal}
Le point $H$ appartient à $d$ et à $\Delta$. Or $\Delta\perp d$. Donc $H$ est le projeté orthogonal de $A$ sur $d$.
\end{correctionbox}
\begin{correctionbox}{6. Distance $AH$}
\[
AH=\sqrt{\left(4-\frac{12}{5}\right)^2+\left(1-\frac95\right)^2}
=\sqrt{\left(\frac85\right)^2+\left(-\frac45\right)^2}
=\sqrt{\frac{80}{25}}
=\frac{4\sqrt5}{5}.
\]
\end{correctionbox}
\begin{correctionbox}{7. Droite parallèle à $d$ passant par $A$}
Une droite parallèle à $d$ a le même vecteur normal, donc une équation de la forme :
\[
2x-y+c=0.
\]
Comme elle passe par $A(4;1)$ :
\[
2\times4-1+c=0,
\qquad c=-7.
\]
Une équation est donc :
\[
2x-y-7=0.
\]
\end{correctionbox}

\subsection*{Sujet 3 -- Cercle, diamètre et tangente}
\begin{correctionbox}{1. Centre du cercle}
Le centre du cercle de diamètre $[AB]$ est le milieu de $[AB]$ :
\[
I\left(\frac{-1+5}{2};\frac{2+4}{2}\right)=(2;3).
\]
\end{correctionbox}
\begin{correctionbox}{2. Rayon}
\[
AB=\sqrt{(5+1)^2+(4-2)^2}=\sqrt{36+4}=2\sqrt{10}.
\]
Le rayon vaut donc :
\[
r=\frac{AB}{2}=\sqrt{10}.
\]
\end{correctionbox}
\begin{correctionbox}{3. Équation du cercle}
Le centre est $I(2;3)$ et le rayon est $\sqrt{10}$, donc :
\[
(x-2)^2+(y-3)^2=10.
\]
\end{correctionbox}
\begin{correctionbox}{4. Critère du diamètre}
On a :
\[
\overrightarrow{MA}=(-1-x;2-y),\qquad
\overrightarrow{MB}=(5-x;4-y).
\]
Donc :
\[
\overrightarrow{MA}\cdot\overrightarrow{MB}=(-1-x)(5-x)+(2-y)(4-y).
\]
En développant :
\[
\overrightarrow{MA}\cdot\overrightarrow{MB}=x^2+y^2-4x-6y+3.
\]
Or :
\[
(x-2)^2+(y-3)^2-10=x^2+y^2-4x-6y+3.
\]
Ainsi :
\[
\overrightarrow{MA}\cdot\overrightarrow{MB}=0
\iff
(x-2)^2+(y-3)^2=10.
\]
Donc $M\in\mathcal C$ si et seulement si $\overrightarrow{MA}\cdot\overrightarrow{MB}=0$.
\end{correctionbox}
\begin{correctionbox}{5. Appartenance de $T$}
Pour $T(1;0)$ :
\[
(1-2)^2+(0-3)^2=1+9=10.
\]
Donc $T$ appartient au cercle.
\end{correctionbox}
\begin{correctionbox}{6. Vecteur normal à la tangente}
La tangente au cercle en $T$ est perpendiculaire au rayon $[IT]$. Donc un vecteur normal à la tangente est :
\[
\overrightarrow{IT}=(1-2;0-3)=(-1;-3).
\]
\end{correctionbox}
\begin{correctionbox}{7. Équation de la tangente}
La tangente a pour vecteur normal $(-1;-3)$ et passe par $T(1;0)$. Elle a donc une équation de la forme :
\[
-1(x-1)-3(y-0)=0.
\]
Donc :
\[
-x+1-3y=0,
\qquad x+3y-1=0.
\]
Une équation de la tangente est :
\[
x+3y-1=0.
\]
\end{correctionbox}
\end{document}
